% ============================================================
% 高考数学满分突破 · 不等式的统一视角
% 编译命令：xelatex -interaction=nonstopmode 文件名.tex
% ============================================================
\documentclass[12pt, a4paper]{ctexart}
\usepackage[margin=2.2cm]{geometry}
\usepackage{amsmath, amssymb}
\usepackage{enumitem}
\usepackage{tcolorbox}
\usepackage{booktabs}
\usepackage{array}
\usepackage{fancyhdr}
\usepackage{xcolor}
\usepackage{tikz}
\usepackage{hyperref}
\usepackage{graphicx}
\usepackage{multirow}
\usepackage{pgfplots}
\pgfplotsset{compat=1.18}

% 颜色定义
\definecolor{mainblue}{RGB}{41, 98, 168}
\definecolor{lightblue}{RGB}{230, 240, 255}
\definecolor{tipgreen}{RGB}{34, 139, 87}
\definecolor{lightgreen}{RGB}{235, 250, 240}
\definecolor{warnred}{RGB}{200, 50, 50}
\definecolor{lightred}{RGB}{255, 240, 240}
\definecolor{examorange}{RGB}{210, 120, 20}
\definecolor{lightorange}{RGB}{255, 248, 235}
\definecolor{apppurple}{RGB}{120, 60, 170}
\definecolor{lightpurple}{RGB}{245, 238, 255}

% 页眉页脚
\pagestyle{fancy}
\fancyhf{}
\fancyhead[L]{\small\textcolor{mainblue}{\textbf{高考数学满分突破 · 不等式的统一视角}}}
\fancyhead[R]{\small\textcolor{mainblue}{适用：高三冲刺}}
\fancyfoot[C]{\small —\ \thepage\ —}
\renewcommand{\headrulewidth}{0.6pt}
\renewcommand{\headrule}{\hbox to\headwidth{\color{mainblue}\leaders\hrule height \headrulewidth\hfill}}
\renewcommand{\footrulewidth}{0pt}

% tcolorbox 样式
\newtcolorbox{knowledgebox}[1][]{%
  colback=lightblue, colframe=mainblue,
  fonttitle=\bfseries\large, title=#1,
  arc=3mm, boxrule=0.8pt,
  left=8pt, right=8pt, top=6pt, bottom=6pt
}

\newtcolorbox{tipbox}[1][]{%
  colback=lightgreen, colframe=tipgreen,
  fonttitle=\bfseries, title=#1,
  arc=2mm, boxrule=0.6pt,
  left=6pt, right=6pt, top=4pt, bottom=4pt
}

\newtcolorbox{warnbox}[1][]{%
  colback=lightred, colframe=warnred,
  fonttitle=\bfseries, title=#1,
  arc=2mm, boxrule=0.6pt,
  left=6pt, right=6pt, top=4pt, bottom=4pt
}

\newtcolorbox{examplebox}[1][]{%
  colback=lightorange, colframe=examorange,
  fonttitle=\bfseries, title=#1,
  arc=3mm, boxrule=0.8pt,
  left=8pt, right=8pt, top=6pt, bottom=6pt
}

\newtcolorbox{appbox}[1][]{%
  colback=lightpurple, colframe=apppurple,
  fonttitle=\bfseries, title=#1,
  arc=3mm, boxrule=0.8pt,
  left=8pt, right=8pt, top=6pt, bottom=6pt
}

\setCJKmainfont{Hiragino Sans GB}
\setCJKsansfont{Heiti SC}

\begin{document}

% ============ 标题页 ============
\begin{center}
{\Huge\bfseries\textcolor{mainblue}{高考数学满分突破}}\\[8pt]
{\LARGE\bfseries\textcolor{mainblue}{不等式的统一视角}}\\[12pt]
\rule{0.6\textwidth}{1pt}\\[8pt]
{\normalsize 适用：高三冲刺\quad |\quad 课时：2课时}
\end{center}
\vspace{10pt}

% ============ 第一部分：知识点讲解 ============
\section*{\textcolor{mainblue}{一、知识点讲解}}

\subsection*{1.1\quad 约束优化与拉格朗日乘数法直观}

\begin{knowledgebox}[约束优化问题的一般形式]
在约束条件 $g(x, y) = 0$ 下，求目标函数 $f(x, y)$ 的最大值或最小值。
这类问题是高考中最值问题的统一本质。
\end{knowledgebox}

\begin{knowledgebox}[拉格朗日乘数法的几何直觉]
构造拉格朗日函数 $L(x, y, \lambda) = f(x, y) - \lambda \cdot g(x, y)$。

在最优解处，目标函数 $f$ 的梯度与约束 $g$ 的梯度\textbf{平行}：
\[ \nabla f = \lambda \nabla g \]
\textbf{几何意义}：最优点处，目标函数的等值线与约束曲线相切。
\end{knowledgebox}

\begin{tipbox}[高考中的直观应用]
虽然高考不要求写出拉格朗日乘数法的完整形式，但其思想贯穿高考最值问题：
\begin{enumerate}[label=(\arabic*)]
  \item 求 $f(x,y)$ 在 $g(x,y) = 0$ 下的最值 $\Longleftrightarrow$ 等值线与约束相切
  \item 线性目标函数在可行域上的最值 $\Longleftrightarrow$ 目标函数等值线平移到可行域边界
  \item 基本不等式的取等条件 $\Longleftrightarrow$ 拉格朗日条件的特殊情况
\end{enumerate}
\end{tipbox}

\begin{warnbox}[易错提醒]
拉格朗日乘数法给出的是\textbf{必要条件}而非充分条件，找到的点可能是最大值、最小值或鞍点。在高考中，通常需要结合几何意义或边界检验来确认。
\end{warnbox}

\subsection*{1.2\quad 线性规划的几何本质}

\begin{knowledgebox}[线性规划问题]
\textbf{目标函数}：$z = ax + by$（线性）

\textbf{约束条件}：
\begin{itemize}
  \item 若干线性不等式 $\begin{cases} a_1 x + b_1 y \le c_1 \\ a_2 x + b_2 y \le c_2 \\ \vdots \end{cases}$
\end{itemize}

\textbf{可行域}：所有约束的交集形成的凸多边形区域。

\textbf{最优解}：目标函数 $z$ 在可行域上的最大值或最小值。
\end{knowledgebox}

\begin{tipbox}[几何本质——等值线平移法]
目标函数 $z = ax + by$ 可改写为 $y = -\dfrac{a}{b}x + \dfrac{z}{b}$（$b \ne 0$）。

这是一族斜率为 $-\dfrac{a}{b}$ 的平行直线。将直线平移：
\begin{enumerate}[label=(\arabic*)]
  \item 若 $b > 0$，$z$ 越大则截距越大，直线向"上"平移
  \item 若 $b < 0$，$z$ 越大则截距越小，直线向"下"平移
\end{enumerate}
最优解出现在直线恰好经过可行域的某个\textbf{顶点}时。
\end{tipbox}

\begin{warnbox}[线性规划的核心结论]
\textbf{定理}：若线性规划问题有最优解，则最优解必在可行域的某个顶点处取得。

因此，只需计算所有顶点处的目标函数值，比较即可。
\end{warnbox}

\subsection*{1.3\quad 优化思想看高考最值}

\begin{knowledgebox}[高考最值问题的统一框架]
高考中的最值问题，本质上都是在某种约束下求某个目标函数的极值。统一视角如下：

\begin{center}
\begin{tabular}{lll}
\toprule
\textbf{题型} & \textbf{目标函数} & \textbf{约束条件} \\
\midrule
线性规划 & $z = ax + by$ & 线性不等式组 \\
基本不等式 & $f(x,y)$（如 $xy$, $x+y$） & $g(x,y) = c$（常数） \\
三角函数最值 & $a\sin\theta + b\cos\theta$ & $\theta \in \mathbb{R}$ \\
导数求最值 & $f(x)$ & $x \in [a, b]$ 或 $f'(x) = 0$ \\
解析几何最值 & 距离/面积 & 曲线方程 \\
\bottomrule
\end{tabular}
\end{center}
\end{knowledgebox}

\begin{tipbox}[基本不等式的统一推导]
由拉格朗日乘数法视角：在 $x + y = s$（常数）下求 $xy$ 的最大值。

$L = xy - \lambda(x + y - s)$。令 $\dfrac{\partial L}{\partial x} = y - \lambda = 0$，$\dfrac{\partial L}{\partial y} = x - \lambda = 0$。

得 $x = y = \dfrac{s}{2}$，即 $xy \le \left(\dfrac{s}{2}\right)^2 = \dfrac{s^2}{4}$。

这就是 $xy \le \left(\dfrac{x+y}{2}\right)^2$，即 AM--GM 不等式的特殊情况。
\end{tipbox}

\begin{knowledgebox}[常用不等式工具箱]
\begin{enumerate}[label=(\arabic*)]
  \item \textbf{均值不等式链}（$a, b > 0$）：
  \[ \frac{2}{\frac{1}{a}+\frac{1}{b}} \le \sqrt{ab} \le \frac{a+b}{2} \le \sqrt{\frac{a^2+b^2}{2}} \]
  即调和平均 $\le$ 几何平均 $\le$ 算术平均 $\le$ 平方平均。

  \item \textbf{柯西--施瓦茨不等式}：
  \[ (a_1^2 + a_2^2)(b_1^2 + b_2^2) \ge (a_1 b_1 + a_2 b_2)^2 \]
  等号条件：$\dfrac{a_1}{b_1} = \dfrac{a_2}{b_2}$。

  \item \textbf{权方和不等式}：
  \[ \frac{a^2}{x} + \frac{b^2}{y} \ge \frac{(a+b)^2}{x+y} \quad (x, y > 0) \]
\end{enumerate}
\end{knowledgebox}

% ============ 第二部分：例题精讲 ============
\section*{\textcolor{examorange}{二、例题精讲}}

\begin{examplebox}[【例题1】基本不等式——条件最值]
已知 $x > 0$，$y > 0$，且 $x + y = 4$，求 $\dfrac{1}{x} + \dfrac{4}{y}$ 的最小值。
\end{examplebox}

\textbf{解：}利用权方和不等式（或柯西--施瓦茨不等式）：
\[ \frac{1}{x} + \frac{4}{y} = \frac{1^2}{x} + \frac{2^2}{y} \ge \frac{(1+2)^2}{x+y} = \frac{9}{4} \]
等号条件：$\dfrac{1}{x} = \dfrac{2}{y}$，即 $y = 2x$。结合 $x + y = 4$，得 $x = \dfrac{4}{3}$，$y = \dfrac{8}{3}$。

最小值为 $\boxed{\dfrac{9}{4}}$。

\begin{examplebox}[【例题2】线性规划——等值线法]
设 $x, y$ 满足约束条件 $\begin{cases} x + y \le 4 \\ x - y \le 2 \\ x \ge 0,\; y \ge 0 \end{cases}$，求 $z = 2x + 3y$ 的最大值。
\end{examplebox}

\textbf{解：}画出可行域，顶点为 $O(0, 0)$，$A(2, 0)$，$B(3, 1)$，$C(0, 4)$。

逐一计算：
\begin{itemize}
  \item $z_O = 0$
  \item $z_A = 2 \times 2 + 3 \times 0 = 4$
  \item $z_B = 2 \times 3 + 3 \times 1 = 9$
  \item $z_C = 2 \times 0 + 3 \times 4 = 12$
\end{itemize}

\textbf{等等}，验证 $C(0, 4)$ 是否在可行域内：$0 + 4 = 4 \le 4$，$0 - 4 = -4 \le 2$，满足。

最大值为 $\boxed{12}$，在 $(0, 4)$ 处取得。

\begin{examplebox}[【例题3】拉格朗日乘数法——二元条件最值]
已知 $x^2 + y^2 = 1$，求 $3x + 4y$ 的最大值。
\end{examplebox}

\textbf{解：}\textbf{方法一（柯西不等式）}：
\[ (3x + 4y)^2 \le (3^2 + 4^2)(x^2 + y^2) = 25 \times 1 = 25 \]
故 $3x + 4y \le 5$，等号条件 $\dfrac{x}{3} = \dfrac{y}{4}$，即 $x = \dfrac{3}{5}$，$y = \dfrac{4}{5}$。

\textbf{方法二（拉格朗日乘数法）}：

$L = 3x + 4y - \lambda(x^2 + y^2 - 1)$。

$\dfrac{\partial L}{\partial x} = 3 - 2\lambda x = 0$，$\dfrac{\partial L}{\partial y} = 4 - 2\lambda y = 0$。

得 $x = \dfrac{3}{2\lambda}$，$y = \dfrac{4}{2\lambda}$。代入 $x^2 + y^2 = 1$：

$\dfrac{9 + 16}{4\lambda^2} = 1$，$\lambda^2 = \dfrac{25}{4}$，$\lambda = \pm\dfrac{5}{2}$。

当 $\lambda = \dfrac{5}{2}$ 时，$x = \dfrac{3}{5}$，$y = \dfrac{4}{5}$，$3x + 4y = 5$。

最大值为 $\boxed{5}$。

\begin{examplebox}[【例题4】三角函数最值]
求函数 $f(\theta) = \sin\theta + \sqrt{3}\cos\theta$ 的最大值和最小值。
\end{examplebox}

\textbf{解：}$f(\theta) = \sin\theta + \sqrt{3}\cos\theta = 2\left(\dfrac{1}{2}\sin\theta + \dfrac{\sqrt{3}}{2}\cos\theta\right) = 2\sin\!\left(\theta + \dfrac{\pi}{3}\right)$。

由于 $-1 \le \sin\!\left(\theta + \dfrac{\pi}{3}\right) \le 1$，故最大值为 $\boxed{2}$，最小值为 $\boxed{-2}$。

\begin{tipbox}[辅助角公式的优化视角]
$a\sin\theta + b\cos\theta = \sqrt{a^2+b^2}\sin(\theta + \varphi)$。

本质上是在约束 $\sin^2\theta + \cos^2\theta = 1$ 下求 $a\sin\theta + b\cos\theta$ 的最值，与例题3 同构。
\end{tipbox}

\begin{examplebox}[【例题5】导数法求最值]
求函数 $f(x) = x^3 - 3x + 2$ 在 $[-3, 3]$ 上的最大值和最小值。
\end{examplebox}

\textbf{解：}$f'(x) = 3x^2 - 3 = 3(x-1)(x+1)$。

令 $f'(x) = 0$，得 $x = -1$ 或 $x = 1$。

计算函数值：
\begin{itemize}
  \item $f(-3) = -27 + 9 + 2 = -16$
  \item $f(-1) = -1 + 3 + 2 = 4$
  \item $f(1) = 1 - 3 + 2 = 0$
  \item $f(3) = 27 - 9 + 2 = 20$
\end{itemize}

最大值为 $\boxed{20}$（$x = 3$），最小值为 $\boxed{-16}$（$x = -3$）。

\begin{examplebox}[【例题6】柯西不等式——巧解最值]
已知 $a + b = 1$，$a, b > 0$，求 $\sqrt{2a + 1} + \sqrt{2b + 1}$ 的最大值。
\end{examplebox}

\textbf{解：}由柯西--施瓦茨不等式：
\[ \left(\sqrt{2a+1} + \sqrt{2b+1}\right)^2 \le (1^2 + 1^2)\left((2a+1) + (2b+1)\right) \]
\[ = 2 \times (2(a+b) + 2) = 2 \times 4 = 8 \]
故 $\sqrt{2a+1} + \sqrt{2b+1} \le 2\sqrt{2}$。

等号条件：$2a + 1 = 2b + 1$，即 $a = b = \dfrac{1}{2}$。

最大值为 $\boxed{2\sqrt{2}}$。

\begin{examplebox}[【例题7】解析几何中的最值]
已知椭圆 $\dfrac{x^2}{4} + y^2 = 1$，$P$ 为椭圆上的点，求 $P$ 到直线 $x + y - 4 = 0$ 的最小距离。
\end{examplebox}

\textbf{解：}设 $P(2\cos\theta, \sin\theta)$，则 $P$ 到直线的距离为
\[ d = \frac{|2\cos\theta + \sin\theta - 4|}{\sqrt{2}} \]
由于 $2\cos\theta + \sin\theta \le \sqrt{4+1} = \sqrt{5} < 4$，故 $2\cos\theta + \sin\theta - 4 < 0$。

\[ d = \frac{4 - (2\cos\theta + \sin\theta)}{\sqrt{2}} \]
$d$ 最小当 $2\cos\theta + \sin\theta$ 最大，即 $\sqrt{5}$。

\[ d_{\min} = \frac{4 - \sqrt{5}}{\sqrt{2}} = \frac{(4 - \sqrt{5})\sqrt{2}}{2} = \boxed{2\sqrt{2} - \frac{\sqrt{10}}{2}} \]

\begin{examplebox}[【例题8】不等式证明——AM-GM]
已知 $a, b, c > 0$，$abc = 1$，证明：$a + b + c \ge 3$。
\end{examplebox}

\textbf{证：}由 AM-GM 不等式：
\[ \frac{a + b + c}{3} \ge \sqrt[3]{abc} = \sqrt[3]{1} = 1 \]
故 $a + b + c \ge 3$。等号条件：$a = b = c = 1$。

\begin{examplebox}[【例题9】含参不等式——分类讨论]
已知 $a > 0$，$b > 0$，求 $\dfrac{a}{b} + \dfrac{b}{a}$ 的最小值。
\end{examplebox}

\textbf{解：}设 $t = \dfrac{a}{b} > 0$，则 $f(t) = t + \dfrac{1}{t}$。

$f'(t) = 1 - \dfrac{1}{t^2}$，令 $f'(t) = 0$ 得 $t = 1$（$t > 0$）。

$f''(1) = 2 > 0$，故 $t = 1$ 为极小值点，$f(1) = 2$。

\textbf{或直接用 AM-GM}：$t + \dfrac{1}{t} \ge 2\sqrt{t \cdot \dfrac{1}{t}} = 2$，等号当 $t = 1$ 即 $a = b$。

最小值为 $\boxed{2}$。

\begin{examplebox}[【例题10】综合——不等式与数列]
设等差数列 $\{a_n\}$ 的前 $n$ 项和为 $S_n$，已知 $S_7 = 49$，$S_{17} = 289$，求 $S_n$ 的最小值。
\end{examplebox}

\textbf{解：}$S_n = na_1 + \dfrac{n(n-1)}{2}d$。

由 $S_7 = 7a_1 + 21d = 49$，得 $a_1 + 3d = 7$。

由 $S_{17} = 17a_1 + 136d = 289$，得 $a_1 + 8d = 17$。

两式相减：$5d = 10$，$d = 2$，$a_1 = 1$。

$S_n = n + n(n-1) = n^2$。

故 $S_n = n^2$，最小值在 $n = 1$ 时取得，$S_1 = \boxed{1}$。

\begin{tipbox}[不等式问题的统一策略]
\begin{enumerate}[label=(\arabic*)]
  \item \textbf{识别问题类型}：目标函数是什么？约束条件是什么？
  \item \textbf{选择工具}：基本不等式（条件简单）、柯西不等式（含平方和）、导数（一元）、线性规划（线性约束）
  \item \textbf{验证取等条件}：最值一定在取等条件可达时取得
\end{enumerate}
\end{tipbox}

% ============ 第三部分：应用题 ============
\section*{\textcolor{apppurple}{三、应用题}}

\begin{appbox}[【应用1】最优设计方案]
某工厂用甲、乙两种原料生产 A、B 两种产品。生产一件 A 产品需甲原料 2 kg、乙原料 1 kg，利润 3 万元；生产一件 B 产品需甲原料 1 kg、乙原料 2 kg，利润 4 万元。现有甲原料 100 kg、乙原料 80 kg。问如何安排生产使利润最大？
\end{appbox}

\textbf{解：}设生产 A 产品 $x$ 件，B 产品 $y$ 件。

约束条件：$\begin{cases} 2x + y \le 100 \\ x + 2y \le 80 \\ x \ge 0,\; y \ge 0 \end{cases}$

目标函数：$z = 3x + 4y$。

可行域顶点：$(0, 0)$，$(50, 0)$，$(0, 40)$，交点由 $2x + y = 100$ 和 $x + 2y = 80$ 解得 $(40, 20)$。

计算：
\begin{itemize}
  \item $(0, 0)$：$z = 0$
  \item $(50, 0)$：$z = 150$
  \item $(0, 40)$：$z = 160$
  \item $(40, 20)$：$z = 120 + 80 = 200$
\end{itemize}

最大利润为 $\boxed{200}$ 万元，生产 A 产品 40 件，B 产品 20 件。

\begin{appbox}[【应用2】用料最优化]
一个容积为 $V$ 的圆柱形罐头，问底面半径 $r$ 和高 $h$ 满足什么关系时，用料（表面积）最小？
\end{appbox}

\textbf{解：}约束：$\pi r^2 h = V$，即 $h = \dfrac{V}{\pi r^2}$。

表面积：$S = 2\pi r^2 + 2\pi rh = 2\pi r^2 + \dfrac{2V}{r}$。

求导：$S'(r) = 4\pi r - \dfrac{2V}{r^2}$，令 $S'(r) = 0$：

$4\pi r = \dfrac{2V}{r^2}$，$r^3 = \dfrac{V}{2\pi}$，$r = \sqrt[3]{\dfrac{V}{2\pi}}$。

此时 $h = \dfrac{V}{\pi r^2} = \dfrac{V}{\pi} \cdot \left(\dfrac{2\pi}{V}\right)^{2/3} = 2\sqrt[3]{\dfrac{V}{2\pi}} = 2r$。

用料最小时，底面直径等于高，即 $h = 2r$。

\begin{appbox}[【应用3】投资收益最大化]
某投资者有资金 10 万元，可投资两个项目。项目 A 的收益率为 $20\%$，但投资额不能超过 6 万元；项目 B 的收益率为 $10\%$，无上限。若总收益要求不低于 1.5 万元，求投资方案中项目 B 的最大投资额。
\end{appbox}

\textbf{解：}设投资 A 项目 $x$ 万元，B 项目 $y$ 万元。

约束条件：$\begin{cases} x + y = 10 \\ 0 \le x \le 6 \\ 0.2x + 0.1y \ge 1.5 \end{cases}$

由 $x + y = 10$ 得 $y = 10 - x$，代入收益约束：$0.2x + 0.1(10-x) \ge 1.5$，即 $0.1x + 1 \ge 1.5$，$x \ge 5$。

结合 $x \le 6$：$5 \le x \le 6$。故 $y = 10 - x \in [4, 5]$。

B 项目最大投资额为 $\boxed{5}$ 万元（此时 $x = 5$）。

% ============ 第四部分：课堂练习 ============
\section*{\textcolor{tipgreen}{四、课堂练习}}

\begin{enumerate}
  \item 已知 $x > 0$，$y > 0$，$2x + 3y = 6$，求 $xy$ 的最大值。
  \item 已知 $x, y$ 满足 $\begin{cases} x + y \le 6 \\ x \ge 0 \\ y \ge 0 \end{cases}$，求 $z = x - 2y$ 的最大值和最小值。
  \item 已知 $x^2 + y^2 = 4$，求 $x + y$ 的最大值。
  \item 求函数 $f(x) = x + \dfrac{4}{x}$（$x > 0$）的最小值。
  \item 已知 $a, b > 0$，$a + 2b = 1$，求 $\dfrac{1}{a} + \dfrac{1}{b}$ 的最小值。
  \item 已知正数 $x, y$ 满足 $x + 3y = 5xy$，求 $3x + 4y$ 的最小值。
\end{enumerate}

\begin{knowledgebox}[参考答案]
\begin{enumerate}
  \item 由 AM-GM：$\dfrac{2x + 3y}{2} \ge \sqrt{2x \cdot 3y} = \sqrt{6xy}$。$2x + 3y = 6$，$3 \ge \sqrt{6xy}$，$xy \le \dfrac{3}{2}$。
  等号当 $2x = 3y = 3$，$x = \dfrac{3}{2}$，$y = 1$。$xy$ 最大值为 $\boxed{\dfrac{3}{2}}$。

  \item 可行域为以 $(0,0)$、$(6,0)$、$(0,6)$ 为顶点的三角形。
  $z_O = 0$，$z_{(6,0)} = 6$，$z_{(0,6)} = -12$。
  最大值 $\boxed{6}$，最小值 $\boxed{-12}$。

  \item 由柯西不等式：$(x+y)^2 \le 2(x^2+y^2) = 8$，$x+y \le 2\sqrt{2}$。
  等号当 $x = y = \sqrt{2}$。最大值为 $\boxed{2\sqrt{2}}$。

  \item AM-GM：$x + \dfrac{4}{x} \ge 2\sqrt{x \cdot \dfrac{4}{x}} = 4$。等号当 $x = 2$。最小值为 $\boxed{4}$。

  \item 权方和：$\dfrac{1}{a} + \dfrac{1}{b} = \dfrac{1^2}{a} + \dfrac{1^2}{b} \ge \dfrac{(1+1)^2}{a+b}$。但需用 $a + 2b = 1$：
  $\dfrac{1}{a} + \dfrac{1}{b} = \left(\dfrac{1}{a} + \dfrac{1}{b}\right)(a + 2b) = 1 + \dfrac{2b}{a} + \dfrac{a}{b} + 2 = 3 + \dfrac{2b}{a} + \dfrac{a}{b}$。
  由 AM-GM：$\dfrac{2b}{a} + \dfrac{a}{b} \ge 2\sqrt{2}$。
  最小值为 $\boxed{3 + 2\sqrt{2}}$。

  \item 由 $x + 3y = 5xy$ 得 $\dfrac{1}{y} + \dfrac{3}{x} = 5$。
  $3x + 4y = (3x + 4y) \cdot \dfrac{5}{5} = \dfrac{1}{5}(3x + 4y)\left(\dfrac{3}{x} + \dfrac{1}{y}\right)$。
  展开：$\dfrac{1}{5}\left(9 + \dfrac{3x}{y} + \dfrac{12y}{x} + 4\right) = \dfrac{1}{5}\left(13 + \dfrac{3x}{y} + \dfrac{12y}{x}\right) \ge \dfrac{1}{5}(13 + 2\sqrt{36}) = \dfrac{13 + 12}{5} = 5$。
  等号当 $\dfrac{3x}{y} = \dfrac{12y}{x}$，即 $x = 2y$。代入 $\dfrac{3}{2y} + \dfrac{3}{x} = 5$ 验证。
  最小值为 $\boxed{5}$。
\end{enumerate}
\end{knowledgebox}

\vfill
\begin{center}
\small\textcolor{gray}{本讲义仅供教学参考，例题可根据学生水平调整难度。}
\end{center}

\end{document}
